%  ============================================================================
% CHAPTER 3: RESULTS
% ============================================================================

\section{Results}

\subsection{Basic Money Distribution}

In this scenario, we consider a simple model of money distribution among agents without debt. 
The total money of the system is conserved, and no agent can have a negative balance.
The agents engage with each other randomly but always with a fixed amount of money in each transaction.

Said so, we can assume that the money distribution will converge to the Boltzmann-Gibbs distribution as predicted by the theoretical framework.
It is possible to see the results in Figure \ref{fig:NoDebtImage}, where the histogram of the money distribution is plotted along with the theoretical Boltzmann-Gibbs distribution.

\begin{figure}[H]
    \centering
    \includegraphics[width=0.4\textwidth]{NoDebtImage.png}
    \caption{The histogram shows the money distribution of agents after a large number of transactions, while the red line represents the theoretical Boltzmann-Gibbs distribution.
                The simulation was performed with $N=3500$, $M = 10 000$, and $\Delta m = 1$. The simulation was run for $10^6$ transactions.}
    \label{fig:NoDebtImage}
\end{figure}

Later on it will be discussed how the richer agents don't really follow Boltzmann-Gibbs distribution, but they follow a power law distribution, as observed in real economies.
These agents are not in equilibrium with the rest of the system, therefore they are not part of the thermalized system.

\subsection{Money Distribution with Debt}


\subsection{Power Law Distribution}


